% Section 4. Sources: paper2/revised_algorithm.md (4.1-4.7), search_soundness.md,
% solver_fix.md (cost, re-certification), prior_art_counterexample.md (prior
% art), certificates.md (proof object), dataset.md (the dataset). Lean:
% lean/MOSPFormalization/Search/. Tables from tables/ (make_tables.py).
% Page numbers for Chu & Stuckey (2009) are those of the preprint
% literature/chu_stuckey_2009.pdf; for Chu (2011) the printed ones.
\section{An exact open-stacks search, read as a pathwidth solver}
\label{sec:search}

The best exact method for MOSP on the standard benchmarks is the \emph{customer
search} of \citet{ChuStuckey2009}. It searches over the orders in which the
customers' stacks close, not over product sequences, and prunes with dominance
rules, three of them stated as theorems (their Theorems~1, 2 and~3) and a fourth,
the subset rule, in their \S2. \citet{ChuStuckey2009} place MOSP among ``graph
path-width and gate matrix layout'', citing \citet{LY2002}, but do not use
pathwidth in the search. By Theorem~\ref{thm:mosp} it is nevertheless a
pathwidth algorithm: any graph is the MOSP graph of the instance with one
customer per vertex and one product per edge, whose optimum is pathwidth plus
one.

This section reads the search that way, and finds three things. Two of the
published theorems are false, and so are their restatements in Chu's thesis
\citep{Chu2011} and, with a stronger premise, in \citet{Fink2012}. The repaired
first rule is an instance of a known pathwidth theorem, Tamaki's commitment
lemma; the published rule is that lemma with its interior condition dropped. And
the whole repaired search is proved sound in Lean. All Lean names in this
section are in \leanfile{lean/MOSPFormalization/Search/}, which has no
\texttt{sorry} and the standard axioms only.

\subsection{The search decides pathwidth}
\label{sec:search-def}

\citet[p.~4]{ChuStuckey2009} observe that the products matter only through the
MOSP graph. So fix a graph $G=(C,E)$ on the customers and an integer $k\ge0$.
For a customer $c$ and a set $T\subseteq C$ of closed customers, write
$N[c]$ for the closed neighbourhood of $c$, and
\begin{align*}
O(T) &= \textstyle\bigcup_{c\in T}N[c] && \text{the stacks opened once $T$ has closed},\\
\fin(X) &= \{c : N[c]\subseteq X\} && \text{the customers free to close when $X$ is open},\\
\cl(T) &= \fin(O(T)) && \text{the \emph{free closure} of $T$},\\
b(T) &= |O(T)\setminus T| && \text{the stacks open after $T$ has closed},\\
o(c,T) &= N[c]\setminus O(T),\quad \open(c,T)=|o(c,T)| && \text{the stacks closing $c$ would newly open.}
\end{align*}
A \emph{move} at $T$ closes one customer $c\notin T$, at cost
$\cost(T,c)=|O(T\cup\{c\})\setminus T| = b(T\cup\{c\})+1$, the stacks open at
that moment counting $c$; it is \emph{playable} when $\cost(T,c)\le k$. Write
$P_k(T)$ when some order of $C\setminus T$ has every step playable. On entering
a node the search closes every free customer, so it stands only at
\emph{states}, sets with $S=\cl(S)$, and the \emph{child} of $S$ by $c$ is
$S\cdot c=\cl(S\cup\{c\})$. The search's predicate $\Sol_k(S)$ holds when $S=C$,
or when some playable $c$ has $\Sol_k(S\cdot c)$.\lean{stepCost, Solvable, SearchSol}
Every node the search visits satisfies the invariant $b(S)\le k$.

\begin{lemma}[free moves never hurt]\label{lem:F}
If $b(T)\le k$, then $P_k(T)\iff\Sol_k(\cl(T))$.\lean{solvable_iff_searchSol_cl}
\end{lemma}
Closing a customer whose neighbourhood is open opens nothing, never makes a
later step dearer, and leaves one fewer stack open. The hypothesis is needed:
closing the free customers one by one costs $b(T)$ at each step.

\begin{theorem}[the search decides pathwidth]\label{thm:decides}
For nonempty $C$, $\Sol_k(\emptyset)\iff\pw(G)+1\le k$; on the MOSP graph of an
instance $M$ with a requirement, $\Sol_k(\emptyset)\iff Z(M)\le k$.\lean{searchSol_empty_iff_pathwidth_add_one_le, searchSol_mospGraph_iff_mospValue_le}
\end{theorem}
\begin{proof}
For a closing order $c_1,\dots,c_n$ with prefixes $S_i$, the $i$-th step costs
$1+|\partial S_i|$, where $\partial X$ is the set of vertices outside $X$ with a
neighbour in $X$. Read backwards as a layout, the cut separating $C\setminus S_i$
from $S_i$ counts exactly $\partial S_i$. So the cost of the order is the vertex
separation of the reversed layout plus one, and Lemma~\ref{lem:F} at the root and
Theorem~\ref{thm:vs} finish the proof. The Lean statement goes through
narrowness \citep{LY2002ref11}, which counts the forward order directly.
\end{proof}

So a refutation, $\neg\Sol_k(\emptyset)$, may only mean $\pw(G)>k-1$. The rest
of the section asks whether the pruning rules preserve that meaning.

\paragraph{What a rule must satisfy.} At a state $S$ the search has a set $Q$ of
\emph{old moves} (Section~\ref{sec:other-rules}). Its candidates are the
customers outside $S\cup Q$; the \emph{cost cut} keeps the playable ones, $P$,
in index order, and a \emph{filter} keeps a sublist $L\subseteq P$. The filter
is \emph{node-sound} at $(S,Q)$ if, whenever every $q\in Q$ has
$\neg\Sol_k(S\cdot q)$, $\Sol_k(S)$ implies $\Sol_k(S\cdot c)$ for some $c\in L$.\lean{NodeSoundAt}
Each rule justifies a discarded move $r$ by a \emph{covering}, a move $c$ with
$\Sol_k(S\cdot r)\Rightarrow\Sol_k(S\cdot c)$. A node is sound when every link a
rule cites is a true covering and every chain of links from a discarded
candidate ends in $L$ or in $Q$; the second condition can fail through a cycle
even when every link is true. The filter runs the definite move, then the subset
rule, then the better move, each once.

\subsection{The definite move is false as published}
\label{sec:definite}

\citet[\S3.1, p.~6]{ChuStuckey2009} state:
\begin{quote}
\textbf{Theorem 1.} Suppose S ++ [q] is playable and close(q, S) $\ge$
open(q, S), then if U$'$ = S ++ R is a solution, there exists a solution
U = S ++ [q] ++ R$'$,
\end{quote}
where $\close(q,S)$ counts the customers $d$ with $o(d,S)\subseteq o(q,S)$
(p.~4). The rule keeps $q$ alone. We read $d$ as ranging over customers not yet
closed, as the code does; read literally, closed customers count too, which
weakens the premise, and both readings are refuted below. \citet[Thm~6.3.6,
p.~142]{Chu2011} restates the theorem with the same premise and proof, reworded
for $k$ stacks. Write $X=S\cdot q$.

\begin{lemma}[what the premise says]\label{lem:premise}
$\close(q,S)=|X\setminus S|$. Hence the premise holds iff $b(X)\le b(S)$: the
child has no more open stacks than the parent.\lean{closeCount_eq, isDefinite_iff}
\end{lemma}

\begin{figure}[t]
\centering
\includegraphics[width=0.85\linewidth]{definite_counterexample.pdf}
\caption{The graph of Counterexample~\ref{cex:definite} at the state where the
definite move fails. Customer 2 is closed; 1, 3 and 4 are open. Moving 0 opens
the new stacks 0, 7 and 11 and closes 0, 3 and 4, so open = close = 3 and the
rule keeps 0 alone. But 3 and 4 have the single new stack 0 between them:
closing them first gains two stacks for the price of one, which a solution
starting with 0 cannot match. Regenerate with \texttt{python -m
paper2.counterexample\_figure}, which recomputes every labelled fact.}
\label{fig:cex}
\end{figure}

\begin{counterexample}\label{cex:definite}
Let $G$ be the graph on customers $0,\dots,13$ with the 26 edges
\begin{gather*}
0\text{--}3,\ 0\text{--}4,\ 0\text{--}7,\ 0\text{--}11,\ 1\text{--}2,\ 1\text{--}5,\ 1\text{--}6,\ 2\text{--}3,\ 2\text{--}4,\ 5\text{--}9,\ 5\text{--}11,\ 5\text{--}12,\ 6\text{--}8,\\
6\text{--}13,\ 7\text{--}9,\ 7\text{--}10,\ 7\text{--}13,\ 8\text{--}9,\ 8\text{--}10,\ 8\text{--}11,\ 9\text{--}10,\ 9\text{--}11,\ 10\text{--}11,\ 10\text{--}12,\ 10\text{--}13,\ 12\text{--}13,
\end{gather*}
with $S=\{2\}$, $k=6$ and $q=0$ (Figure~\ref{fig:cex}). Then $O(S)=\{1,2,3,4\}$,
$o(0,S)=\{0,7,11\}$ and $\cost(S,0)=6$, so $0$ is playable; the customers with
$o(d,S)\subseteq\{0,7,11\}$ are $0$, $3$ and $4$, so $\close(0,S)=3\ge\open(0,S)$;
and $0$ comes first in index order, so the rule keeps only $0$. A solution
exists from $S$: the order
\[1,3,4,6,12,13,0,5,7,8,9,10,11\]
costs at most 6. None exists
from the child $X=\{0,2,3,4\}$: only seven sets are reachable from $X$ by moves
of cost at most 6, and none is $C$.\lean{definiteMove_counterexample, cexFamily_closed}
\end{counterexample}

The published implication, as a universal statement, is refuted in Lean under
both readings of $\close$, and so is the thesis's Theorem~6.3.6.\lean{chuStuckey_theorem1_false, chuStuckey_theorem1_false_literal, chuThesis_theorem636_false}
The smallest graph on which the rule loses the last solution has 14 customers:
there is none on any graph to 9 vertices, by exhaustion, and none on the 88,802
connected graphs on 10 vertices with at most 14 edges. The counterexample was
found by a search over a family of gadgets built from the failed proof, and
minimised by deletion. \citet[pp.~27--28]{Fink2012} restates the rule with a
stronger premise, counting only dominated customers of smaller index than $q$.
The graph above refutes it under no labelling, but a third twin of 3 and 4 gives
a 15-customer graph that does.\lean{fink_theorem1_false}

\paragraph{Why the published proof fails.} The proof moves $q$ to the front of a
solution and claims that after each prefix ``we have at most open(q, S) extra
stacks open, but at least close(q, S) extra stacks closed'' (p.~6). At a prefix
$T\supseteq S$ that does not yet contain $q$, moving $q$ forward changes the open
stacks by $b(T\cup X)-b(T) = |N[q]\setminus O(T)|-|X\setminus T|$. At $T=S$ this
is $\open-\close\le0$. But customers of $X$ that the solution has \emph{already}
closed are in $T$ and no longer count as extra stacks closed, while the stacks
$q$ opens may still be new. In the example, a solution may begin with 3, after
which 4 is free: at $T=\{2,3,4\}$, moving 0 forward costs one more stack. A
second, independent slip: even where the conclusion holds, the literal sequence
$S \mathbin{+\!\!+} [q] \mathbin{+\!\!+} R'$ need not be a solution, because a
customer that $q$ frees stays open until its turn; the argument needs the free
moves.

\paragraph{The repair.} Call $q$ \emph{hereditarily definite} at $S$ when
$b(X)\le b(B)$ for every $B$ with $S\subseteq B\subseteq X$ and $q\notin B$. By
Lemma~\ref{lem:premise}, the published premise is the case $B=S$. In the
counterexample the repair fails at $B=\{2,3,4\}$, where $b(B)=2<3=b(X)$.\lean{IsHereditarilyDefinite, not_isHereditarilyDefinite_cex}

\begin{lemma}[open stacks are submodular]\label{lem:submod}
For all $A,B\subseteq C$, $b(A\cup B)+b(A\cap B)\le b(A)+b(B)$.\lean{openStacks_submodular}
\end{lemma}
\begin{proof}
$b(T)=|O(T)|-|T|$, $O(A\cup B)=O(A)\cup O(B)$ and $O(A\cap B)\subseteq
O(A)\cap O(B)$; the customers themselves balance exactly.
\end{proof}

\begin{theorem}[the repaired definite move is sound]\label{thm:repaired}
If $q\notin S$ is hereditarily definite at $S$, then $P_k(S)\Rightarrow P_k(X)$.
At a state with the invariant, $\Sol_k(S)\Rightarrow\Sol_k(S\cdot q)$.\lean{solvable_cl_insert_of_hereditarilyDefinite, searchSol_cl_insert_of_hereditarilyDefinite}
\end{theorem}
\begin{proof}
Take an order from $S$ of cost at most $k$ with prefixes $T_i$, in which $q$ is
the $j$-th step, and follow it from $X$, skipping customers already in $X$. A
step $i<j$ closes $c_i$ at $T_{i-1}\cup X$. With $Z=T_i$ we have
$S\subseteq Z\cap X\subseteq X$ and $q\notin Z$, so the hypothesis gives
$b(X)\le b(Z\cap X)$, and Lemma~\ref{lem:submod} gives $b(Z\cup X)\le b(Z)$: the
step costs no more than the original. After step $j-1$ the closed set contains
$T_j$ with the same opened set, and the rest of the order applies. The search's
form follows by Lemma~\ref{lem:F}.
\end{proof}

This theorem is not new. The hereditary premise is the commitment condition of
Tamaki, and the proof is the proof of the commitment lemma,
\citet[Lemma~1, p.~142]{Kitsunai2016}; Section~\ref{sec:layouts} makes the
correspondence exact. What we did not find in the papers we hold is the
following cheap form of the test.

\begin{theorem}[the repair is a matching condition]\label{thm:matching}
Let $Y=X\setminus(S\cup\{q\})$, and join each $d\in Y$ to the stacks of
$o(d,S)\subseteq o(q,S)$. Then $q$ is hereditarily definite at $S$ iff this
bipartite graph has a matching with at least $\open(q,S)-1$ edges.\lean{isHereditarilyDefinite_iff_hasDefiniteMatching}
\end{theorem}
\begin{proof}
The sets of the repair are $B=S\cup D$, $D\subseteq Y$, and
$b(X)\le b(S\cup D)$ iff $|D|-|o(D)|\le\delta$ with
$\delta=|Y|+1-\open(q,S)$, where $o(D)$ is the union of the $o(d,S)$. By the
deficiency form of Hall's theorem the largest matching has size
$|Y|-\max_D(|D|-|o(D)|)$, which is at least $|Y|-\delta=\open(q,S)-1$ exactly
when the repair holds. The Lean proof reduces the deficiency form to Mathlib's
Hall theorem with $\delta$ dummy stacks.
\end{proof}

In words: the rule is safe exactly when all but one of the new stacks can be
paired with a different freed customer who needed it. In the counterexample,
closing 0 opens three new stacks and frees 3 and 4, both of whom needed only
stack 0: one pairing where two are required. Every $q$ with $\open(q,S)\le1$
passes at once, so the published rule is right there.\lean{isHereditarilyDefinite_of_openCount_le_one}

\subsection{The other rules}
\label{sec:other-rules}

\paragraph{The subset rule.} ``If $o(c_i,S)\subseteq o(c_j,S)$ and $i<j$ \dots\
move $j$ can be removed'' \citep[\S2, p.~5]{ChuStuckey2009}. As coded, $d$
dominates $r$ when $o(d,S)\subseteq o(r,S)$ and the inclusion is strict or $d<r$;
every dominated playable $r$ is dropped, with a fallback that keeps $P$ if all
would go. The covering is sound for every set $S$: if $o(d,S)\subseteq o(r,S)$
and $r$ is playable, then $d$ is playable and $\Sol_k(S\cdot r)\Rightarrow
\Sol_k(S\cdot d)$, because $d$ is free after $r$ and closing $d$ then $r$ reaches
the same state at no greater cost. With the index tie-break, domination is a
strict order, so an undominated move with a solution survives, and the fallback
never fires when a solution exists.\lean{searchSol_cl_insert_of_newlyOpened_subset, subsetFilter_sound}
The tie-break is necessary: without it, two customers with equal new stacks
dominate each other and both go, which loses a solution on a 7-customer
graph.\lean{noTieBreak_counterexample}

\paragraph{The better move.} \citet[\S3.2, p.~6]{ChuStuckey2009}: ``Suppose
S ++ [q] and S ++ [r, q] are playable and close(q, S $\cup$ \{r\}) $\ge$
open(q, S $\cup$ \{r\}) then if U$'$ = S ++ [r] ++ R is a solution there exists a
solution U = S ++ [q] ++ R$'$.'' Its premise is Theorem~1's premise for $q$ at the
child $\cl(S\cup\{r\})$, together with playability of $r$ then $q$, and its
conclusion swaps $q$ and $r$.\lean{isBetter_iff} The swap is sound; the appeal to
Theorem~1 is not. On the graph of Counterexample~\ref{cex:definite} at the root
with $k=6$, the definite move does not fire, customers 0 and 2 both survive the
subset rule, and the premise holds for $r=2$, $q=0$, so the rule drops 2 citing
0. But $\{2\}$ has a solution and $\{0\}$ has none.\lean{betterMove_counterexample, chuStuckey_theorem2_false}
At that node another kept customer still has a solution, so this is a false link
and not a lost node; no node where the better move loses the last solution has
been found. \citet[Thm~6.3.8, p.~143]{Chu2011} states the better move with a
different premise, $\close(q,S)\ge\open(q,S\cup\{r\})$, and proves it through
Theorem~6.3.6. The same graph refutes it at $S=\{2\}$, $r=3$, $q=0$.\lean{chuThesis_theorem638_false, chuThesis_theorem638_false_literal}
The repair cites $q$ for $r$ only if $q$ is hereditarily definite at
$\cl(S\cup\{r\})$, one matching at the child; then the swap argument goes
through.\lean{IsRepairedBetter, searchSol_cl_insert_of_repairedBetter}

Two earlier faults of this rule in our own implementation are now proved in
Lean too: a close count that included the customers $r$ finishes on its own,
and an order of the rules in which the subset rule could cite a customer the
better move had already dropped, which loses a node through a cycle of true
links.\lean{bugA_counterexample, bugB_counterexample}
The second is an instance of the general observation of
\citet[Prop.~3, Example~6]{Beck2025} that mutually dominating transitions must
be tie-broken.

\paragraph{The old move.} \citet[Thm~3, p.~7]{ChuStuckey2009} let a child inherit
a sibling $q$ whose subtree failed, provided that inserting $q$ before each
later step keeps that step playable. The rule is sound as coded: if
$\cost(S\cup\{q\},c)\le k$, then
$\Sol_k(\cl(\{q\}\cup S\cdot c))\Rightarrow\Sol_k(S\cdot q)$, so a refuted
sibling stays refuted along any path whose steps pass the test.\lean{searchSol_reinsert, not_searchSol_of_oldMove}
The test is necessary; the smallest graph that shows it is the path on five
vertices.\lean{reinsert_needs_test}

\paragraph{The memo.} A refuted state is recorded and refused when met again.
Whether a state with $b(S)\le k$ has a solution depends on the state alone, not
on the path that reached it.\lean{solvable_iff_of_cl_eq} One implementation of
ours declined to run the memo beside the old move, on the ground that a
refutation obtained with inherited old moves depends on the ancestors. A node
taken on its own, with an old move that is not in fact refuted, can indeed
record a solvable state;\lean{exec_fake_oldMove} the next theorem shows that this
never happens in a real run.

\subsection{The repaired search is sound}
\label{sec:theorem}

\begin{proposition}\label{prop:filter}
At every state with the invariant, for every set $Q$ of refuted old moves, the
filter definite $\to$ subset $\to$ better with both repairs is node-sound.\lean{repairedFullFilter_sound}
\end{proposition}
\begin{proof}
If the definite move fires, Theorem~\ref{thm:repaired} applies. Otherwise the
subset rule keeps a survivor with a solution; let $w$ be the earliest. If the
better move dropped $w$, it cited an earlier survivor that the repaired
argument proves solvable, contradicting the choice of $w$.
\end{proof}

\begin{theorem}[runs are sound]\label{thm:runs}
Let the filter be node-sound and return only playable customers. If at the start
of a run the invariant holds, every memo entry is a genuine refutation, and
every old move is refuted, then the run leaves only genuine memo entries, and if
it answers \emph{false} at $S$ then $\neg\Sol_k(S)$. This holds for any child
order and any memo policy, with the old move on or off.\lean{Exec.sound}
\end{theorem}
The proof is an induction over the run, in the order in which subtrees complete:
every old move a node holds is a sibling subtree that completed with \emph{false}
earlier in the same run, or an inherited one that passed the reinsertion test.

\begin{theorem}[the revised search is sound]\label{thm:main}
Run the customer search with free moves, the memo, old moves and the filter
definite $\to$ subset $\to$ better, the definite move firing only on a
hereditarily definite customer and the better move citing only under the
repaired premise. If it refutes $k$ from the root, then no closing order costs at
most $k$, $\pw(G)>k-1$, and on the MOSP graph of an instance $M$ with a
requirement, $Z(M)>k$.\lean{exec_repairedFullFilter_sound, exec_repairedFullFilter_pathwidth, exec_repairedFullFilter_mospValue}
\end{theorem}

What the Lean model covers is the search as specified; that our C and Python
implement it is checked, not proved. A transcription of the Lean filter agreed
with the production filter on 16,244,090 node checks over graphs on 1 to 17
vertices; the matching agreed with the hereditary premise on 36,939,226
candidates; and a search written from the theorems agreed with the C, the Python
and an exact oracle on 17,431,232 runs over 52,592 graphs on 1 to 17
vertices.\footnote{\texttt{paper2/solver\_fix.md}, items 01--05, and
\texttt{paper2/revised\_algorithm.md}, section~4.4.4, with the regenerate
commands.} Under the published rules the filter is not
node-sound,\lean{not_codeFilterSound_cexGraph} and in a check of 570,206 runs with
memo and old move it lost the last solution at a node in 58 runs, all on graphs
built around Counterexample~\ref{cex:definite} and all at $k$ equal to the
optimum. Every one found a solution through another branch and answered
correctly. No whole instance on which the published rules answer \emph{unsat}
wrongly is known.

\subsection{The search in pathwidth language}
\label{sec:layouts}

For a set $T$ of vertices write $N(T)$ for its \emph{border}, the vertices
outside $T$ with a neighbour inside, and $d(T)=|N(T)|$, in the notation of
\citet[p.~390]{KobayashiKomuroTamaki2014}. A vertex sequence is $w$-feasible when
every prefix has $d\le w$, and $\vs(G)$ is the least $w$ for which some
permutation is $w$-feasible. Table~\ref{tab:dictionary} translates.

\begin{table}[t]
\centering\small
\caption{The customer search in the vocabulary of exact pathwidth algorithms.}
\label{tab:dictionary}
\begin{tabularx}{\linewidth}{@{}>{\raggedright}p{3.6cm}>{\raggedright}X>{\raggedright\arraybackslash}p{5.2cm}@{}}
\toprule
Customer search & Layout & Lean \\
\midrule
closing order from $\emptyset$ & vertex sequence, read forwards & \path{orderOfLayout} \\
open stacks $b(T)$ & border size $d(T)$ & \path{openStacks_eq_card_boundary} \\
$\cost(T,c)$ & $d(T\cup\{c\})+1$ & \path{stepCost_eq_card_boundary_add_one} \\
$\Sol_k(\emptyset)$ & $\vs(G)\le k-1$ & \path{searchSol_empty_iff_pathwidth_add_one_le} \\
free closure $\cl(T)$ & full set of $T$ \citep[p.~328]{SuchanVillanger2009} & \path{mem_cl_iff} \\
memo & failure table & -- \\
definite move, repaired & commitment to the child & \path{isHereditarilyDefinite_iff_isCommittable} \\
\bottomrule
\end{tabularx}
\end{table}

\citet[Algorithm~1, p.~393]{KobayashiKomuroTamaki2014} decide whether a vertex
set extends to a $w$-feasible permutation: replace the set by a committed
extension, answer from a failure table or at $V$, recurse on each feasible
one-vertex extension, record failures. The customer search is this loop with
$w=k-1$: the free closure and the definite move play the part of the commitment,
the memo is the failure table, and the cost cut is the feasibility test.
\citet[\S3.2]{CoudertMazauricNisse2014} has the same three parts in a branch and
bound. Neither line cites the other in the papers we hold: \citet{ChuStuckey2009}
predates the commitments and both 2014 papers, and none of the pathwidth papers
cites it.

\paragraph{Commitments.} \citet[p.~142]{Kitsunai2016}, stating Tamaki's notion,
call an extension $\tau$ of $\sigma$ \emph{committable} when $d(X')\ge d(V(\tau))$
for every $X'$ with $V(\sigma)\subseteq X'\subseteq V(\tau)$; their Lemma~1 says
that if $\sigma$ extends to a feasible permutation, so does $\tau$.

\begin{lemma}\label{lem:endpoint}
For $q\notin S$ and $X=\cl(S\cup\{q\})$: the published premise holds iff
$d(X)\le d(S)$, the commitment condition at the single set $X'=S$; and the
repaired premise holds iff $X$ is a commitment from $S$, the condition at every
$X'$ between them.\lean{isDefinite_iff_endpoint, isHereditarilyDefinite_iff_isCommittable}
\end{lemma}

So Theorem~\ref{thm:repaired} is the commitment lemma at the child, and the
published Theorem~1 is a commitment whose interior condition was dropped. In
Counterexample~\ref{cex:definite} the endpoint test passes with $d(X)=3=d(S)$,
while the interior set $\{2,3,4\}$ has border 2.\lean{cex_isDefinite_not_isCommittable}
The published premise is not wasted: by \citet[Lemma~10, Cor.~2,
pp.~148--149]{Kitsunai2016} it guarantees a commitment to the least-border set
between $S$ and $X$, which need not contain $q$.\lean{exists_isCommittable_of_isDefinite}
In the counterexample that set is $\{2,3,4\}$.\lean{cex_isCommittable_234}
Committing there when the matching fails is a second repair that gives up no
pruning; we have not built it.

A commitment of depth one, $S$ to $S\cup\{q\}$, holds iff $\open(q,S)\le1$.\lean{isCommittable_insert_iff}
The greedy step of \citet[Lemma~3, p.~49]{CoudertMazauricNisse2014} is exactly
that condition,\lean{isGreedyStep_iff} and the push rule of
\citet[Rule~1, p.~330]{SuchanVillanger2009} is a case of it. So every published
pathwidth rule of this kind lies in the region $\open\le1$, where the published
definite move is right; \citet[p.~388, Table~2]{KobayashiKomuroTamaki2014} report
depth-one commitments as ``extremely effective''. Chu and Stuckey's Theorem~1
reaches into $\open\ge2$, and that is where it fails. For the cost of the test,
\citet[Cor.~2]{Kitsunai2016} find the least border between the ends of a given
extension by a minimum $s$--$t$ separator; Theorem~\ref{thm:matching} is that
computation for the single target $X$, where it becomes a bipartite matching.

Table~\ref{tab:known} sums up what is known, limited to the pathwidth papers we
hold \citep{SuchanVillanger2009,Bodlaender2012,CoudertMazauricNisse2014,CoudertMazauricNisse2016,KobayashiKomuroTamaki2014,Kitsunai2016}.
``None found'' is not a claim of priority.

\begin{table}[t]
\centering\small
\caption{Each rule of the customer search against the exact-pathwidth literature.}
\label{tab:known}
\begin{tabularx}{\linewidth}{@{}>{\raggedright}p{3.2cm}>{\raggedright}X>{\raggedright\arraybackslash}p{4.1cm}@{}}
\toprule
Rule & Published counterpart & Status \\
\midrule
free move & full sets \citep{SuchanVillanger2009}; \citet[Prop.~3]{Kitsunai2016} & known \\
definite move, $\open\le1$ & depth-one commitments; greedy step; push rule & known \\
definite move, published & -- & false (Cex.~\ref{cex:definite}) \\
definite move, repaired & commitment lemma \citep[Lemma~1]{Kitsunai2016} & known; matching test (Thm~\ref{thm:matching}) not found \\
subset rule & none found & proved here \\
better move, published & none found & false \\
better move, repaired & none found & proved here \\
old move & none found & proved here \\
memo & subset DP \citep{Bodlaender2012}; prefix and failure tables & known; with the old move, proved here (Thm~\ref{thm:runs}) \\
\bottomrule
\end{tabularx}
\end{table}

\paragraph{Prior reports.} We looked for any earlier report of the error. We
listed the works citing \citet{ChuStuckey2009}, Chu's thesis and its IJCAI 2013
abstract in OpenAlex, Semantic Scholar and OpenCitations: 141 works. Of the 38
that cite the CP paper we read 24 in full (63\%), and of the 103 that cite only
the thesis, 32 (31\%).\footnote{\texttt{python3 -m paper2.prior\_art\_sweep};
the table of citers is in \texttt{paper2/prior\_art\_counterexample.md}.} No
work read gives a counterexample, an erratum or a correction. One,
\citet{Fink2012}, restates the definite move as true with a stronger premise,
refuted above. Three use an implementation of the rules and report no failure;
among them \citet{Frinhani2018} ran the original code to obtain the optima that
\citet{MartinYanassePinto2022} reuse. We found no prior report of the error.
The pathwidth literature has the correct form of the rule, as a commitment, and
none of the pathwidth papers we hold mentions MOSP.

\subsection{The revised algorithm in practice}
\label{sec:practice}

The revised search changes two tests. The definite move fires on the first
playable $q$ that passes the published test and then the matching of
Theorem~\ref{thm:matching}; a $q$ that fails the matching is passed over. The
better move cites $q$ for $r$ only if the same matching succeeds at
$\cl(S\cup\{r\})$. The matching joins at most $\close-1$ customers to $\open$
stacks and stops at $\open-1$ edges. Since 1 October 2026 these are the default
in both of our solvers, the MOSP customer search and the graph pathwidth solver,
in C and in Python; the published rules remain behind a flag.

\begin{table}[t]
\centering\small
\caption{The cost of the repair: nodes under the published and the repaired
rules, in paired runs of the same decision in one worker. Totals over the pairs
that finished under both settings; a censored pair says nothing about cost.
MOSP rows refute the optimum minus one; the graph row is whole descents at
120\,s each.}
\label{tab:cost}
\resizebox{\linewidth}{!}{% Generated by python -m paper2.latex.make_tables. Source: paper2/data/solver_fix_cost_{mosp40,cs,cs125,pw}.csv.
% Do not edit by hand.
\begin{tabular}{@{}llrrrrr@{}}
\toprule
Where & Vertices & Pairs finished & Nodes, published & Nodes, repaired & Ratio & Worst pair \\
\midrule
MOSP corpus, \texttt{default} & 9--40 & 6,135 of 6,135 & 228,147 & 228,154 & 1.00003 & 1.009 \\
MOSP corpus, \texttt{csearch} & 9--40 & 6,135 of 6,135 & 202,962 & 202,971 & 1.00004 & 1.040 \\
Chu \& Stuckey classes & 50--100 & 118 of 125 & 1,358,096,671 & 1,362,205,210 & 1.00303 & 1.013 \\
Chu \& Stuckey classes, cap $2\times10^8$ nodes & 125 & 11 of 23 & 90,969,508 & 90,961,419 & 0.99991 & 1.000 \\
Graph pathwidth, whole descents & 4--2{,}916 & 731 of 880 & 2,977,849,807 & 3,038,670,717 & 1.020 & -- \\
\bottomrule
\end{tabular}
}
\end{table}

Table~\ref{tab:cost} gives the cost in nodes. The repair changes tree sizes by
at most a few percent and never changed an answer. Per node it costs about 2\% in
the MOSP C engine, measured on six hard instances at 75--125 customers. The
published test passes and the matching then fails on 0.07--0.27\% of
definite-move candidates, a rate that rises with size.

\paragraph{Re-certification.} Of the 6,374 certified optima in our MOSP corpus,
6,259 do not depend on the customer search: they rest on an exact subset-lattice
computation, a checked DRAT refutation, an earlier SAT search, or a lower bound
equal to the value. The other 115 rested on the customer search under the
published rules. The repaired solver has re-refuted the value minus one on 112
of them, and no value changed. Table~\ref{tab:split} gives the seven hardest,
125-customer instances at densities 2 and 4, which needed a split of the root
over many cores; three are still running at the time of writing, and keep their
values as verified upper bounds until they finish.\footnote{\texttt{python -m
paper2.solver\_fix\_split -{}-summary}; the split is proved sound in
\leanfile{Search/Split.lean}.}

\begin{table}[t]
\centering\small
\caption{Re-certification of the hardest corpus values under the repaired rules,
by a split of the root into independent tasks. \emph{$k$}: the value minus one,
which must be refuted. For an instance in flight, the nodes and core-hours are
those of the finished tasks so far.}
\label{tab:split}
% Generated by python -m paper2.latex.make_tables. Source: paper2/data/solver_fix_split_{tasks,results}.csv, read 2026-10-03 09:01.
% Do not edit by hand.
\begin{tabular}{@{}lrrlrrr@{}}
\toprule
Instance & Value & $k$ & State & Tasks refuted & Nodes & Core-hours \\
\midrule
\texttt{Random-125-125-4-1\_0} & 57 & 56 & refuted & 669 & $4.97\times10^{10}$ & 18.1 \\
\texttt{Random-125-125-2-4\_0} & 24 & 23 & refuted & 1,686 & $4.98\times10^{11}$ & 154.1 \\
\texttt{Random-125-125-4-5\_0} & 46 & 45 & refuted & 616 & $7.96\times10^{10}$ & 26.1 \\
\texttt{Random-125-125-4-2\_0} & 57 & 56 & refuted & 842 & $7.64\times10^{10}$ & 30.2 \\
\texttt{Random-125-125-2-1\_0} & 24 & 23 & in flight & 197 & $8.24\times10^{10}$ & 51.0 \\
\texttt{Random-125-125-2-5\_0} & 20 & 19 & in flight & 168 & $6.97\times10^{10}$ & 45.6 \\
\texttt{Random-125-125-4-4\_0} & 51 & 50 & in flight & 745 & $1.28\times10^{11}$ & 67.1 \\
\bottomrule
\end{tabular}
% read 2026-10-03 09:01; last task stamp 2026-10-03T09:00:59

\end{table}

\paragraph{Certificates.} A refutation by the search is a claim about a tree of
$10^{11}$ nodes at 125 customers, and nothing in Theorem~\ref{thm:main} lets a
reader check that the code ran the search it models. So the search emits a
certificate: the tree, every pruning step with its rule and witness, and for
each definite and better step the matching of Theorem~\ref{thm:matching}. An
independent checker, which shares no code with the search, recomputes every
premise from the matrix, verifies each matching (distinct customers, distinct
stacks, each stack newly opened by its customer, at least $\open-1$ edges),
requires every chain of covers to end at a child or an old move with no cycle,
and replays the tree. Each accepted link is a Lean theorem; that an acyclic
covering by such links is sound is a short paper argument, not yet formal.
Table~\ref{tab:certs} gives the results on the corpus at 9--75 customers.
Every refutation emitted within 60\,s verifies, 6,275 of 6,286 under the
\texttt{csearch} configuration, and none is rejected; the 35.5\,MB of
certificates check in 147\,s. On the graph of Counterexample~\ref{cex:definite}
rooted at $\{2\}$, the published rules emit an exhaustive tree that claims no
solution; a checker of the published premises accepts it, and ours rejects it,
even with a largest matching attached (one edge where two are needed).

\begin{table}[t]
\centering\small
\caption{Certificates for the repaired rules on the certified corpus at 9--75
customers, refuting the optimum minus one, 60\,s emission budget.
\emph{Unknown}: emission did not finish. Below 41 customers the
\texttt{default} rows match \texttt{csearch} and are omitted.}
\label{tab:certs}
\resizebox{\linewidth}{!}{% Generated by python -m paper2.latex.make_tables. Source: paper2/data/certificates/repaired.csv.
% Do not edit by hand.
\begin{tabular}{@{}llrrrrrrrrr@{}}
\toprule
Customers & Config. & Inst. & Verified & Rejected & Unknown & Definite & Better & Matching edges & gz bytes (median) & Check ms (median) \\
\midrule
9--10 & \texttt{csearch} & 1,614 & 1,614 & 0 & 0 & 307 & 92 & 32 & 332 & 0.04 \\
11--20 & \texttt{csearch} & 2,508 & 2,508 & 0 & 0 & 2,657 & 1,171 & 658 & 356 & 0.08 \\
21--30 & \texttt{csearch} & 1,816 & 1,816 & 0 & 0 & 14,796 & 6,148 & 5,675 & 414 & 0.24 \\
31--40 & \texttt{csearch} & 197 & 197 & 0 & 0 & 46,337 & 18,176 & 17,150 & 650 & 1.27 \\
41--50 & \texttt{csearch} & 120 & 120 & 0 & 0 & 530,424 & 120,302 & 152,597 & 12,045 & 33.24 \\
51--60 & \texttt{csearch} & 1 & 1 & 0 & 0 & 5,155 & 2,709 & 1,885 & 126,508 & 278.86 \\
61--75 & \texttt{csearch} & 30 & 19 & 0 & 11 & 103,095 & 31,404 & 37,261 & 26,516 & 143.70 \\
\midrule
41--50 & \texttt{default} & 120 & 120 & 0 & 0 & 748,952 & 0 & 191,780 & 12,665 & 37.01 \\
51--60 & \texttt{default} & 1 & 1 & 0 & 0 & 24,685 & 0 & 8,797 & 661,951 & 1198.43 \\
61--75 & \texttt{default} & 30 & 20 & 0 & 10 & 312,897 & 0 & 110,118 & 34,162 & 187.03 \\
\bottomrule
\end{tabular}
% total csearch: 6275 of 6286 verified, 0 rejected, 35.47 MB gz, 147.4 s, 3353245 nodes
% total default: 6276 of 6286 verified, 0 rejected, 54.77 MB gz, 242.3 s, 5361282 nodes
}
\end{table}

\subsection{A dataset with certified pathwidth}
\label{sec:dataset}

Section~\ref{sec:equiv} makes one certified pathwidth an exact answer to every
problem in Table~1's exact core: pathwidth, vertex separation and the vertex
separator game take the width $w$; MOSP, gate matrix layout with multiple
folding, one-dimensional logic, interval thickness, narrowness and node search
take $w+1$. We therefore built one dataset from every benchmark collection of
those problems that we could obtain, and the MOSP corpus. The unit is the
isomorphism class of the graph, the MOSP graph for a matrix instance: two
instances with isomorphic graphs have the same value under every problem, so the
dataset holds one record per class, with its members and their isomorphism maps.
Each record carries the graph, the width, its provenance (\emph{refutation},
\emph{bound} or \emph{upper bound}), a witness layout of exactly that width and,
for a bound, a minor certificate that replays.

\begin{table}[t]
\centering\small
\caption{The dataset by collection. \emph{Classes}: isomorphism classes within
the collection (nauty certificates; 78 graphs above 5,000 vertices by file
content only). \emph{Owned}: classes whose first member is in this collection,
so the column sums to the dataset's class count. The provenance columns count
owned classes with a witness layout.}
\label{tab:dataset}
\resizebox{\linewidth}{!}{% Generated by python -m paper2.latex.make_tables. Source: paper2/data/dataset/{index,classes}.csv.gz.
% Do not edit by hand.
\begin{tabular}{@{}lrrrrrrr@{}}
\toprule
Collection & Files & Classes & Owned & Refutation & Bound & Upper bound & No witnessed value \\
\midrule
MOSP: Challenge 2005 & 5,806 & 3,169 & 3,169 & 3,168 & 1 & 0 & 0 \\
MOSP: Challenge 2005 (second copy) & 46 & 46 & 0 & 0 & 0 & 0 & 0 \\
MOSP: Chu \& Stuckey & 200 & 200 & 200 & 198 & 0 & 2 & 0 \\
MOSP: Faggioli \& Bentivoglio & 300 & 287 & 274 & 274 & 0 & 0 & 0 \\
MOSP: SCOOP & 24 & 24 & 24 & 24 & 0 & 0 & 0 \\
VLSI gate matrix circuits & 11 & 11 & 11 & 5 & 5 & 1 & 0 \\
VSPLIB trees & 50 & 28 & 28 & 28 & 0 & 0 & 0 \\
VSPLIB grids & 50 & 50 & 50 & 9 & 0 & 19 & 22 \\
VSPLIB Harwell-Boeing & 73 & 72 & 72 & 32 & 7 & 33 & 0 \\
Small (Marti et al. 2008) & 84 & 84 & 84 & 32 & 52 & 0 & 0 \\
Rome graphs & 11,534 & 11,199 & 11,199 & 8,126 & 2,746 & 326 & 1 \\
TreewidthLIB colouring & 58 & 58 & 58 & 18 & 13 & 27 & 0 \\
freetdi named graphs & 150 & 146 & 143 & 97 & 22 & 23 & 1 \\
freetdi control-flow graphs & 1,817 & 1,070 & 1,064 & 398 & 661 & 4 & 1 \\
MOSP: Frinhani et al. 2018, large & 610 & 610 & 610 & 0 & 0 & 0 & 610 \\
PACE 2016 & 291 & 288 & 112 & 37 & 18 & 51 & 6 \\
PACE 2017 exact & 200 & 199 & 193 & 61 & 7 & 115 & 10 \\
PACE 2017 heuristic & 200 & 197 & 173 & 11 & 3 & 45 & 114 \\
PACE 2017 bonus & 100 & 100 & 100 & 23 & 0 & 77 & 0 \\
MOSP: Carvalho \& Soma 2015 & 150 & 150 & 150 & 11 & 0 & 139 & 0 \\
\midrule
All & 21,754 &  & 17,714 & 12,552 & 3,535 & 862 & 765 \\
\bottomrule
\end{tabular}
}
\end{table}

Table~\ref{tab:dataset} gives the counts. Twenty collections, 21,754 instance
files, reduce to 17,714 classes: the MOSP collections
\citep{SmithGent2005,FaggioliBentivoglio1998,ChuStuckey2009,CarvalhoSoma2015,Frinhani2018},
eleven VLSI gate matrix circuits \citep{OliveiraLorena2002}, the VSPLIB sets
\citep{Duarte2012} and the Small graphs \citep{MartiCamposPinana2008}, the Rome
graphs \citep{DiBattista1997}, and treewidth collections used as pathwidth
benchmarks \citep{TreewidthLIB,freetdi,PACE2016,PACE2017}. Of the classes, 16,087
carry a certified pathwidth with a witness layout and 862 a witnessed upper
bound. An independent checker, sharing no code with the solvers, re-reads all
20,987 member files with its own readers, maps each onto its record by
isomorphism, re-measures every witness and replays every bound certificate, and
reports no failure on the 16,949 records.\footnote{\texttt{python
paper2/dataset\_check.py -{}-sources}; builder \texttt{python -m paper2.dataset}.}

Against published values: all 84 Small graphs agree with the pathwidths of
\citet[Tables~4--5]{Mallach2018}, and ten of the eleven VLSI circuits are
certified at their published best-known track counts \citep[Table~I]{OliveiraLorena2002};
the eleventh, W4 with 202 nets, is open at 28 against 27. One caution about MOSP
benchmark files emerged: the same public repository ships the Carvalho and Soma
instances as customers by patterns and the Chu and Stuckey instances beside them
as patterns by customers, and only published densities and values tell them
apart. Read the wrong way round, an instance certifies at a different value.

What is left is priced. One more budget step for every open class that fits our
engine (1,024 vertices per component) costs at most 500 core-hours, and on the
Rome graphs each such step has closed about 40\% of what it attempted. The
largest MOSP instances, 400 to 1,000 customers, and the 154 classes above 1,024
vertices are beyond an exact search. Refutation-certified records rest on the
repaired search; for those within reach of the certificate emitter, the proof
object above is the route by which a third party can check them without our code.
